% GENERATED FILE. Edit canonical lessons, then run npm run generate:books.
% canonical-source-hash: fnv1a64:5cc8e262aee3d449
% canonical-lessons: PT-C298-veredicto, PT-C298-recorrer, PT-C298-recurso, PT-C298-recluso, PT-C298-libertar

\chapter{Depois do veredicto}
\label{ch:pt-depois-do-veredicto}
\hlchaptermodality{298}

\begin{chapteropening}
\textbf{By the end of this chapter:} \emph{I can talk about a verdict, an appeal, and a prisoner being set free.}

Complete the last lesson of chapter 298: I can talk about a verdict, an appeal, and a prisoner being set free.
\end{chapteropening}

\section[o veredicto]{o veredicto — the verdict}

\label{lesson:PT-C298-veredicto}



\begin{quote}
Before the new one: say the Portuguese for to sentence, then the Portuguese for to acquit.
\end{quote}

\subsection*{You'll want to know: o veredicto}
\textbf{o veredicto} — \textquotedblleft{}the verdict\textquotedblright{}.

\begin{quote}
O veredicto saiu ao fim da tarde.
\end{quote}

\begin{quote}
\emph{The verdict came out late in the afternoon.}
\end{quote}

The decision of a court on the facts.

\begin{cousinweb}
\emph{veredicto} goes back to Latin \emph{vērē dictum} \textquotedblleft{}truly said\textquotedblright{}, from \emph{vērus} \textquotedblleft{}true\textquotedblright{} and \emph{dīcere} \textquotedblleft{}to say\textquotedblright{} (Lewis \& Short), probably through English \emph{verdict}.
\end{cousinweb}

\begin{grammarlens}[title={What you've built}]
One more word for telling, explaining and describing at B1.
\end{grammarlens}

\subsection*{Your turn}
\begin{itemize}
\raggedright
  \item \emph{Say it:} \emph{o veredicto}
  \item \emph{Say it:} \emph{o veredicto}, once more
  \item \emph{Say it:} a sentence of your own with \emph{o veredicto}
  \item \emph{Recall:} say the Portuguese for to sentence, then the Portuguese for to acquit, then say \emph{o veredicto} again
\end{itemize}

\begin{tcolorbox}[breakable,colback=teal!4,colframe=teal!35!black,title={Before you move on}]
What does o veredicto mean? (\textquotedblleft{}The verdict\textquotedblright{}.) Say it once more.
\end{tcolorbox}
\section[recorrer]{recorrer — to appeal (in court); to turn to (for help)}

\label{lesson:PT-C298-recorrer}



\begin{quote}
Before the new one: say the Portuguese for to acquit, then the Portuguese for the verdict.
\end{quote}

\subsection*{You'll want to know: recorrer}
\textbf{recorrer} — \textquotedblleft{}to appeal (in court); to turn to (for help)\textquotedblright{}.

\begin{quote}
O advogado vai recorrer da decisão.
\end{quote}

\begin{quote}
\emph{The lawyer is going to appeal against the decision.}
\end{quote}

\emph{Recorrer a alguém} — to turn to someone for help.

\begin{cousinweb}
\emph{recorrer} comes from Latin \emph{recurrere} \textquotedblleft{}to run back, to have recourse\textquotedblright{} (EtymDB), built on \emph{currere} \textquotedblleft{}to run\textquotedblright{}.
\end{cousinweb}

\begin{grammarlens}[title={What you've built}]
One more word for telling, explaining and describing at B1.
\end{grammarlens}

\subsection*{Your turn}
\begin{itemize}
\raggedright
  \item \emph{Say it:} \emph{recorrer}
  \item \emph{Say it:} \emph{recorrer}, once more
  \item \emph{Say it:} a sentence of your own with \emph{recorrer}
  \item \emph{Recall:} say the Portuguese for to acquit, then the Portuguese for the verdict, then say \emph{recorrer} again
\end{itemize}

\begin{tcolorbox}[breakable,colback=teal!4,colframe=teal!35!black,title={Before you move on}]
What does recorrer mean? (\textquotedblleft{}To appeal (in court); to turn to (for help)\textquotedblright{}.) Say it once more.
\end{tcolorbox}
\section[o recurso]{o recurso — the appeal; a resource}

\label{lesson:PT-C298-recurso}



\begin{quote}
Before the new one: say the Portuguese for the verdict, then the Portuguese for to appeal.
\end{quote}

\subsection*{You'll want to know: o recurso}
\textbf{o recurso} — \textquotedblleft{}the appeal; a resource\textquotedblright{}.

\begin{quote}
O tribunal aceitou o recurso.
\end{quote}

\begin{quote}
\emph{The court accepted the appeal.}
\end{quote}

Also a resource: \emph{os recursos naturais}.

\begin{cousinweb}
\emph{recurso} is a learned borrowing from Latin \emph{recursus} \textquotedblleft{}a running back, a return\textquotedblright{}, from \emph{recurrere} \textquotedblleft{}to run back\textquotedblright{} (Lewis \& Short).
\end{cousinweb}

\begin{grammarlens}[title={What you've built}]
One more word for telling, explaining and describing at B1.
\end{grammarlens}

\subsection*{Your turn}
\begin{itemize}
\raggedright
  \item \emph{Say it:} \emph{o recurso}
  \item \emph{Say it:} \emph{o recurso}, once more
  \item \emph{Say it:} a sentence of your own with \emph{o recurso}
  \item \emph{Recall:} say the Portuguese for the verdict, then the Portuguese for to appeal, then say \emph{o recurso} again
\end{itemize}

\begin{tcolorbox}[breakable,colback=teal!4,colframe=teal!35!black,title={Before you move on}]
What does o recurso mean? (\textquotedblleft{}The appeal; a resource\textquotedblright{}.) Say it once more.
\end{tcolorbox}
\section[o recluso]{o recluso — the inmate, the prisoner (formal)}

\label{lesson:PT-C298-recluso}



\begin{quote}
Before the new one: say the Portuguese for to appeal, then the Portuguese for the appeal.
\end{quote}

\subsection*{You'll want to know: o recluso}
\textbf{o recluso} — \textquotedblleft{}the inmate, the prisoner (formal)\textquotedblright{}.

\begin{quote}
O recluso trabalha na biblioteca da cadeia.
\end{quote}

\begin{quote}
\emph{The inmate works in the prison library.}
\end{quote}

Formal and official; in speech, \emph{o preso}.

\begin{cousinweb}
\emph{recluso} is a learned borrowing from Late Latin \emph{reclūsus} \textquotedblleft{}shut up, enclosed\textquotedblright{}, from \emph{reclūdere}, built on \emph{claudere} \textquotedblleft{}to shut\textquotedblright{} (Lewis \& Short).
\end{cousinweb}

\begin{grammarlens}[title={What you've built}]
One more word for telling, explaining and describing at B1.
\end{grammarlens}

\subsection*{Your turn}
\begin{itemize}
\raggedright
  \item \emph{Say it:} \emph{o recluso}
  \item \emph{Say it:} \emph{o recluso}, once more
  \item \emph{Say it:} a sentence of your own with \emph{o recluso}
  \item \emph{Recall:} say the Portuguese for to appeal, then the Portuguese for the appeal, then say \emph{o recluso} again
\end{itemize}

\begin{tcolorbox}[breakable,colback=teal!4,colframe=teal!35!black,title={Before you move on}]
What does o recluso mean? (\textquotedblleft{}The inmate, the prisoner (formal)\textquotedblright{}.) Say it once more.
\end{tcolorbox}
\section[libertar]{libertar — to set free, to release}

\label{lesson:PT-C298-libertar}



\begin{quote}
Before the new one: say the Portuguese for the appeal, then the Portuguese for the inmate.
\end{quote}

\subsection*{You'll want to know: libertar}
\textbf{libertar} — \textquotedblleft{}to set free, to release\textquotedblright{}.

\begin{quote}
Libertaram o recluso ao fim de dez anos.
\end{quote}

\begin{quote}
\emph{They released the inmate after ten years.}
\end{quote}

\emph{Libertar um refém} — to free a hostage.

\begin{cousinweb}
\emph{libertar} was borrowed from Medieval Latin \emph{lībertāre} (EtymDB), built on \emph{lībertās} \textquotedblleft{}freedom\textquotedblright{}, from \emph{līber} \textquotedblleft{}free\textquotedblright{}.
\end{cousinweb}

\begin{grammarlens}[title={What you've built}]
One more word for telling, explaining and describing at B1.
\end{grammarlens}

\subsection*{Your turn}
\begin{itemize}
\raggedright
  \item \emph{Say it:} \emph{libertar}
  \item \emph{Say it:} \emph{libertar}, once more
  \item \emph{Say it:} a sentence of your own with \emph{libertar}
  \item \emph{Recall:} say the Portuguese for the appeal, then the Portuguese for the inmate, then say \emph{libertar} again
\end{itemize}

\begin{tcolorbox}[breakable,colback=teal!4,colframe=teal!35!black,title={Before you move on}]
What does libertar mean? (\textquotedblleft{}To set free, to release\textquotedblright{}.) Say it once more.
\end{tcolorbox}
